\section{Related Work and Deep Learning Literature Review}
\label{sec:lit_review}

The computational modeling of chemical property and toxicity prediction has evolved through distinct paradigms of molecular representation, progressing from handcrafted discrete descriptors to end-to-end graph representation learning.

\subsection{Molecular Fingerprints and Representation Limits}
Early computational toxicology often used structural keys and circular hashed fingerprints. Duvenaud et al.~\cite{duvenaud2015convolutional} presented a graph neural network whose learned molecular fingerprints generalize the fixed circular-fingerprint construction: their differentiable graph operations replace a fixed fingerprint-computation step and learn task-dependent features. Kearnes et al.~\cite{kearnes2016molecular} described molecular graph convolutions that operate on atom and pair features, and evaluated them against fingerprint methods on ligand-based virtual-screening tasks. Wu et al.~\cite{wu2018moleculenet} introduced MoleculeNet as a benchmark spanning several molecular machine-learning tasks, with common datasets, metrics, and baselines; its proposed evaluation protocols include scaffold splits based on Bemis--Murcko frameworks~\cite{bemis1996properties}.

\subsection{Message Passing Neural Networks on Molecular Graphs}
Gilmer et al.~\cite{gilmer2017neural} expressed several graph neural network variants in the general Message Passing Neural Network (MPNN) framework and explored new variants for molecular property prediction. Their experiments on QM9 show that message functions using bond information and recurrent atom-state updates can be effective, but do not establish that every MPNN must use these choices. Kipf and Welling~\cite{kipf2017gcn} proposed the Graph Convolutional Network (GCN), using a first-order approximation of localized spectral graph convolutions. A standard GCN uses adjacency and node features rather than explicit bond-feature vectors; molecular implementations that omit bond features therefore do not directly represent bond attributes such as bond order or conjugation.

Yang et al.~\cite{yang2019analyzing} evaluated a directed message-passing neural network (D-MPNN, used in Chemprop) that propagates messages along directed bonds, preventing an immediate message from simply returning along the same bond. In parallel, Xiong et al.~\cite{xiong2020pushing} introduced AttentiveFP, using attention over atom and bond representations to form molecular representations; their experiments include toxicity prediction and illustrate attention-based interpretations of molecular substructures. These results do not establish that attention reliably identifies toxicophores in general.

\subsection{Permutation Invariance, Readout Mechanisms, and Regularization}
Molecules lack a canonical linear atom ordering; hence, mapping variable-sized node feature sets $\{\mathbf{h}_v\}_{v \in \mathcal{V}}$ to a fixed-dimensional molecular embedding $\mathbf{z}$ requires a permutation-invariant readout operation. Vinyals et al.~\cite{vinyals2015order} introduced Set2Set, an attention-based recurrent method for producing a permutation-invariant representation of a set. Gilmer et al.~\cite{gilmer2017neural} used Set2Set as one graph readout and reported strong results for it in their QM9 experiments; this does not imply that it will outperform sum or mean pooling in every setting. Cho et al.~\cite{cho2014learning} introduced the Gated Recurrent Unit (GRU) for recurrent sequence models. Its gates control how information is updated and retained; using a GRU to update graph-node states is an architectural choice made in later graph models.

Deep neural networks trained on molecular datasets of this size can overfit, motivating regularization. Srivastava et al.~\cite{srivastava2014dropout} introduced dropout as a method to reduce co-adaptation of feature detectors. Rong et al.~\cite{rong2020dropedge} proposed DropEdge, which randomly removes graph edges during training as data augmentation and to reduce message propagation; their analysis and experiments concern mitigating overfitting and oversmoothing in graph-convolutional networks. Loshchilov and Hutter~\cite{loshchilov2019decoupled} proposed decoupled weight decay (AdamW), separating weight decay from the adaptive gradient update.

\subsection{Handcrafted Hybrids and Honest Empirical Benchmarking}
A useful caution in deep learning for drug discovery is that a more expressive architecture does not automatically outperform classical methods. Yang et al.~\cite{yang2019analyzing} evaluated adding global molecular features, including RDKit descriptors, to learned representations and reported improvements on some of their benchmark comparisons. This supports testing complementary global descriptors, but does not show that every descriptor set improves every graph model or that message passing cannot learn relevant global properties.

Jiang et al.~\cite{jiang2021could} compared graph-based models with descriptor-based machine-learning methods across public molecular datasets. Their results show that relative performance depends on the dataset and model setup: descriptor-based methods can match or outperform graph models on some tasks, while graph models perform better on others. This motivates evaluating both model families under a common protocol; the paper does not prescribe paired significance tests or negative controls, which are methodological choices in the present study.
